\scanpage{368}\mspage{284}
a) Pour que $S$ soit accessible, il f. et il
suffit qu'il existe un cardinal \inserted{infini} $\pi$\deleted{\unsure{, tel que} \unsure{($\alpha$)|(a)},
utile pour $M$,} et une partie
$S_{(\pi)}$ de $M_\pi$, \deleted{stable \unsure{dans}} tels que
$S$ soit formé des objets de $M$ qui
sont isomorphes à un objet de la
forme $\underset{\unsure{j|i|v}\in I}{\underrightarrow{\lim}}\,x_i$, avec les $x_i$ dans $S_\pi$,
et $I$ grd devant $\pi$. En d'autres
termes, il f. et s. qu'il existe \deleted{un
$\pi$ \inserted{$\geq\pi_0$} utile pour $M$,}\inserted{un cardinal \unsure{infini} $\pi$,} et une \deleted{\unsure{part.}} ss-catégorie
pleine $\underline S(\pi)$ de $M_\pi$, tels que
$\underline S$ soit l'image essentielle du
foncteur canonique \deleted{\unsure{2}} \unsure{pl fid|pl. fid.|pl fidèle}
\[
 \mathrm{Ind}_\pi(\underline S(\pi))\longrightarrow M
 \bigl(\xleftarrow{\approx}\mathrm{Ind}_\pi(M_\pi)\bigr)
 \quad\inserted{\text{(si }\pi\text{ utile pour }M\text{)}}\tag*{(4.13.2.1)}
\]
prolongeant l'inclusion canonique
$\underline S(\pi)\unsure{\hookrightarrow|\subset} M$ (cf. 4.12.3).

b) \inserted{Si $S$ est accessible,} On peut trouver des cardinaux
$\pi'_0,\pi'_1$, tels que $\pi'_1\geq\pi'_0\geq\pi_0$
\deleted{avec $\pi'_0\geq\pi_0$, $\pi'_1\geq\pi_1=\unsure{\operatorname{Sup}(\pi_0,\operatorname{Card\,red} M_{\pi_0})}$,
$\unsure{\pi'_0,\pi'_1\geq\operatorname{Sup}(\pi_0,\pi_1)}$, $\operatorname{Card\,red}M_{\pi_1}$}
\inserted{(où $\pi_0$ est un cardinal donné utile pour $M$)}
% archive sheet number: 367
